\scanpage{71}
\uline{NB} Nous appliquerons ces remarques surtout au
cas où $X'$ est un ouvert de $X$. C'est pour
pouvoir prendre de telles images inverses sans
ennui, qu'il n'aurait pas été commode
de supposer les $X_i$ ou les $X_i^*$ non vides,
ou encore que $\unsure{I|i}\mapsto X_i$ est un \uline{plong\textsuperscript{t}}
d'ens ordonnés $I\hookrightarrow\mathcal P(X)$.

Lorsque $X'\to X$ est une imm. loc\textsuperscript{le} \deleted{\unsure{p}} \inserted{propre}
\deleted{(\unsure{plus gén.}} \inserted{mais \uline{\uuline{pas}}} (ou une imm. ouverte !) alors
la formation des images inverses des
parties \deleted{\unsure{les}} fermées de $X$ commute à la
formation des voisinages tubulaires de
une telle partie dans une autre. Notons
d'ailleurs que pour $i<j$, on aura
(sans hyp. d'ailleurs que $X'-X$ soit une
imm. local)
\[
(4)\quad\left\{\begin{aligned}
R'_{ij}&=R_{ij}\times_X X'\\
S'_{ij}&=S_{ij}\times_X X'
\end{aligned}\right.,
\]
% Diagonal marginal note; the local doubts include the crowded attachment of comme.
\margin{Il y a ici une propriété formelle importante des tubes qu'il faut vérifier --- \unsure{ça est valable} pour les topos définis \unsure{comme} germes d'espaces \unsure{en dehors hyp. métriques} (p.ex. $X$ compact --- brrr!)}

d'où, en le cas d'une imm. local \inserted{propre},
des iso
\[
(5)\quad\left\{\begin{aligned}
V'^{\deleted{\unsure{*}}}_{i,j}&\simeq V_{i,j}\times_X X'\\
V'^{*}_{i,j}&\simeq V^*_{i,j}\times_X X'
\end{aligned}\right.\qquad i\leq j
\]

% Thin grey foliation at lower left, in another hand; no dark author page number observed.
\margin{70}
